\section{线性基：参数、调用与改造}

本节紧跟前面的基础模板。先看每个接口传什么，再看哪些不变量不能动，最后看如何替换维护信息或转移。

\subsection{它解决什么问题}

线性基处理的不是普通加法，而是异或意义下的线性空间。典型题面包括：

\begin{itemize}
\tightlist
\item
  从若干数中选子集，使异或和最大/最小；
\item
  判断一个数能否由给定数异或得到；
\item
  在线加入数，查询当前所有数能得到的最大异或值；
\item
  两个集合的异或空间是否相交，或求异或值的第 \(k\) 小。
\end{itemize}

把每个整数看成二进制向量。异或就是在 \(\mathbb F_2\) 上相加，所以每一位只有 \(0/1\)，消元时只需要``最高位主元''。若数据小于 \(2^{60}\)，开 \passthrough{\lstinline!LOG = 61!}；若题目给到 \(2^{100}\)，必须改用 \passthrough{\lstinline!bitset!} 或高精度实现。


\subsection{可直接使用的模板}

\begin{lstlisting}[language={C++}]
template<int LOG = 61>
struct LinearBasis {
    using U = unsigned long long;
    U p[LOG]{}; // p[i] 的最高位是 i

    // 接口：clear()：把全部基向量清零，使当前 LinearBasis 恢复为空线性基；无返回值。
    void clear() { memset(p, 0, sizeof(p)); }

    // 接口：insert(x)：把数值 x 消元后插入线性基；线性相关时保持基不变；会修改当前结构；多测时重新构造或显式清空成员。
    // 参数：待消元并插入当前线性基的整数 x。
    void insert(U x) {
        // 变量：i=线性基插入时从高到低检查的二进制位编号 i。
        for (int i = LOG - 1; i >= 0; --i) {
            if (((x >> i) & 1ULL) == 0) continue;
            if (!p[i]) {
                p[i] = x;
                return;
            }
            x ^= p[i];
        }
    }

    // 接口：can(x)：判断 x 是否能由当前线性基中的向量异或得到；x 消元为 0 时返回 true。
    // 参数：待判断能否被线性基表示的数值 x。
    bool can(U x) const {
        // 变量：i=判断 x 可表示性时从高到低消元的二进制位编号 i。
        for (int i = LOG - 1; i >= 0; --i) {
            if (!((x >> i) & 1ULL)) continue;
            if (!p[i]) return false;
            x ^= p[i];
        }
        return true;
    }

    // 接口：max_xor(x)：返回把 x 与当前线性基中若干向量异或后的最大值；返回类型 U。
    // 参数：求最大异或值时使用的初始值 x。
    U max_xor(U x = 0) const {
        // 变量：i=求最大异或值时从高到低贪心检查的二进制位编号 i。
        for (int i = LOG - 1; i >= 0; --i)
            x = max(x, x ^ p[i]);
        return x;
    }
};
\end{lstlisting}

\passthrough{\lstinline!insert!} 是异或版高斯消元；\passthrough{\lstinline!can!} 是消元判定；\passthrough{\lstinline!max\_xor!} 从高位到低位贪心，若异或后变大就保留。这个顺序不能反过来，否则低位的选择可能破坏高位最优性。


\subsection{例题：P3812 线性基}

题目：\href{https://www.luogu.com.cn/problem/P3812}{P3812 〖模板〗线性基}。题意是给出 \(n\) 个整数，选择任意子集，使它们的异或和最大。

套用步骤：

\begin{enumerate}
\def\labelenumi{\arabic{enumi}.}
\tightlist
\item
  每个输入数都代表一个可以选择或不选择的向量；
\item
  不需要记录``选了哪些数''，因为目标只依赖最终异或值；
\item
  逐个 \passthrough{\lstinline!insert(x)!}；
\item
  输入结束后输出 \passthrough{\lstinline!basis.max\_xor()!}。
\end{enumerate}

例如输入集合为 \passthrough{\lstinline!1, 5, 9, 4!}：插入后可以得到 \passthrough{\lstinline!1 \^ 4 \^ 8 = 13!}，模板从最高位开始尝试，得到最大值 \passthrough{\lstinline!13!}。完整调用如下：

\begin{lstlisting}[language={C++}]
// 接口：main()：读入 n 个整数建立线性基，并输出可取得的最大子集异或值；入口函数；每组测试前按注释清空全局状态。
int main() {
    ios::sync_with_stdio(false);
    cin.tie(nullptr);
    // 变量：n=待插入线性基的整数个数 n。
    int n;
    cin >> n;
    LinearBasis<61> basis;
    while (n--) {
        // 变量：x=当前读入并待插入线性基的整数 x。
        unsigned long long x;
        cin >> x;
        basis.insert(x);
    }
    cout << basis.max_xor() << '\n';
}
\end{lstlisting}

\subsubsection{\texorpdfstring{测试样例：线性基必须按 \texttt{void\ solve()} 包装}{测试样例：线性基必须按 void solve() 包装}}

下面约定每组输入为 \passthrough{\lstinline!n!} 个非负整数，输出任意子集异或和的最大值；第一行是测试组数 \passthrough{\lstinline!T!}。前两组检查零向量、重复向量和多测重建，后两组检查独立主元以及高位贪心。

\begin{lstlisting}
输入
4
1
0
3
5 5 5
4
1 2 4 8
5
8 7 3 10 5

输出
0
5
15
15
\end{lstlisting}

如果第二组大于 \passthrough{\lstinline!5!}，通常是消元时把重复向量当成了新主元；如果第三、四组小于 \passthrough{\lstinline!15!}，优先检查插入方向和 \passthrough{\lstinline!max\_xor!} 是否从最高位向下贪心。

练习升级：\href{https://www.luogu.com.cn/problem/P4301}{P4301 新 Nim 游戏}。它把``异或空间''放进博弈/删数语境中，看到``剩下的数不能异或出 \(0\)''时，应该想到线性基的秩和线性相关，而不是暴力枚举子集。


\subsection{常见错误}

\begin{itemize}
\tightlist
\item
  把线性基当成普通的异或前缀；它维护的是一组向量，插入顺序不影响空间，但主元排列会影响某些``第 \(k\) 小''模板。
\item
  \passthrough{\lstinline!long long!} 的符号位参与比较时容易出错。纯非负数据优先使用 \passthrough{\lstinline!unsigned long long!}。
\item
  求最大值时要从高位到低位；求最小非零值时通常要先把基化成更适合贪心的形式。
\item
  线性基只能描述``每个数至多选一次''的异或组合；若某个数可以无限次使用，问题已经不是这个模板。
\end{itemize}

